\documentclass{article}

\usepackage{amsmath}

\title{Σ-Φ-C Cognitive Collapse Model}
\author{Your Name}
\date{}

\begin{document}

\maketitle

\begin{abstract}
We propose the Σ-Φ-C model, a cognitive dynamical system describing cognition as transitions between Generation (T), Correction (O), and Collapse (∅) states.
\end{abstract}

\section{Introduction}

Cognition is traditionally modeled as parallel processing. We argue it is better described as a state-transition system.

\section{Model Definition}

We define a cognitive system:

\[
C = \langle S, \Delta, I \rangle
\]

where:

\[
S \in \{T, O, \emptyset\}
\]

\section{State Dynamics}

The system evolves as:

\[
S_{t+1} \sim P(S | S_t, I_t)
\]

\section{Collapse Mechanism}

Collapse occurs when:

\[
p_c = \sigma(I(S_t) - \theta)
\]

where:

\[
\sigma(x) = \frac{1}{1 + e^{-x}}
\]

\section{Reconfiguration Phase}

When collapse occurs:

\[
S_t \rightarrow \emptyset
\]

and:

\[
P(S_{t+1} | \emptyset) = \text{Uniform}(T, O)
\]

\section{Conclusion}

The Σ-Φ-C model provides a unified dynamical framework for cognition.

\end{document}
